\paragraph{From ternary encoding to the marked flag.}
\label{inc-ampl-bandera}
The regions also provide the codes $w_\pi,w_e,w_\varphi$ and
$w_{\rm APP}$ in the inherited frame. The map
\[
 \gamma_{A_W}:w\longmapsto(w,wA_W)
\]
preserves the six-component word as its first coordinate and appends its
transform under $A_W$. The relation $A_W^2=-I_6$ fixes the quadratic
compatibility of this realization. The next step,
$c\mapsto\operatorname{supp}(c)$, preserves the nonzero positions. For
weight-six words, the support identifies exactly the pair $\{c,-c\}$, as
proved by the minimum distance of the code. Thus, quotienting by sign produces
the hexads, while the complete code remains available for operations that
require its symbols.

In this chart, the two constructions of the tetrad coincide:
\begin{equation}
 \{i:K_i\geq729\}
 =H_e\cap P_\pi=B_K=\{4,6,9,10\}.
 \label{inc-ampl-cuaterna}
\end{equation}
The left-hand side arises from an integer comparison of the components of
$K$; the right-hand side, from the intersection of the encoded propagation
and closure supports. The threshold belongs to the declared block chart and
remains explicit in the definition of the reader. The full record preserves
the values on which it acts, even when different thresholds produce the same
subset. The equality therefore expresses compatibility between specified
maps.

The support $H^-_\alpha$ completes that tetrad to a hexad. Membership in the
support $P_\pi$, together with the intersection sizes $(4,3,2)$ relative to
$H_e,H_\varphi,H_{\rm APP}$, selects a unique hexad. The proofs that follow
bring together the two determinations: the one obtained from the sign of
$Z_0$ and the one obtained from incidence. The result is organized as
\begin{equation}
 \widetilde{\mathcal I}_*
 =\bigl(B_K\subset H^-_\alpha;
         P_\pi,H_e,H_\varphi,H_{\rm APP}\bigr).
 \label{inc-ampl-marco}
\end{equation}
The inclusion is the flag; the four additional supports constitute its marked
frame. Simultaneous permutations transport the frame and preserve its
inclusions and intersections. The relevant information resides in this
transported configuration and its isomorphism class.

The five regions of $\pi$ recover their internal differentiation here. They
share $w_\pi$ as their visible ternary code, but preserve distinct emissions,
signatures, and multiplicities. In the binary realization on a marked
dodecad, the transverse region corresponds to the octad whose intersection
with the dodecad is the tetrad; the other four regions correspond to the
lifts of the four hexads in its star. The negative region is distinguished by
$H^-_\alpha$, and the chart order individuates the three positive regions.
This correspondence preserves the regional roles and the distribution
$432+4\cdot144=1008$, while also realizing the incidence partition $1+4$.
